% ============================================================================
% CAPÍTULO 4 — IMPLEMENTACIÓN NUMÉRICA
% Objetivo: describir el integrador propio y su validación.
% Resultado que DEBE contener:
%   · Arquitectura del código (spacetimes / shell / integrators / diagnostics).
%   · El RK4 de paso fijo (escrito a mano, por control y por ser formativo).
%   · El paso ADAPTATIVO por step-doubling: estimador de error local (~h^5),
%     controlador h_new = h·S·(tol/err)^{1/(p+1)}, rechazo-y-repite.
%     Motivación: en a>1 el ~99% de la dinámica está en el ~1% final cerca
%     del horizonte; el paso fijo desperdicia pasos en la región lenta.
%   · Validación: benchmark ExactParametric (a<1, a=1) y monitoreo de |C|.
%     OJO: el estimador de error del stepper y |C| son controles INDEPENDIENTES.
% CRITERIO DE ALCANCE: este capítulo justifica hacer o no el integrador a>1.
%   Si tus claims (cap. 5) no necesitan resultados a>1, esto se acota o se
%   reporta como limitación, no como bloqueante.
% ============================================================================
\chapter{Implementación numérica}
\label{cap:numerico}

\section{Arquitectura del simulador}
% TODO: diagrama de módulos y flujo de datos (ShellParams -> Runner -> ...).

\section{Integrador de Runge--Kutta de cuarto orden}
% TODO: esquema RK4; por qué implementación propia (control, formación).

\section{Paso adaptativo por duplicación de paso}
\label{sec:adaptativo}
% TODO: estimador de error local por step-doubling; controlador de paso;
%   tolerancia mixta atol + rtol·|y| por componente.

\section{Validación}
% TODO: comparación contra ExactParametric; drift de |C| (métrica exterior,
%   r>rs); por qué max_C_drift explota cerca de la singularidad (amplificación
%   numérica de b/2R, no error del integrador).
